\section{Lacunas Regulatórias: Navegando em um Vácuo Jurídico}

O compasso de evolução do arcabouço regulatório global é notorialmente dissonante da velocidade da inovação tecnológica. No Brasil, assim como na União Europeia e nos Estados Unidos, os órgãos sanitários encontram-se estruturados para avaliar dispositivos médicos analógicos, insumos farmacológicos ou, no máximo, softwares de processamento de imagem estáticos. Um sistema tão complexo, preditivo e dinâmico como um gêmeo digital odontológico subverte as categorias tradicionais de classificação.

\subsection{A Ambiguidade da Classificação SaMD (Software as a Medical Device)}

A Agência Nacional de Vigilância Sanitária (ANVISA) baseia-se na Resolução RDC 657/2022 para regular softwares em saúde, aderindo aos princípios do \textit{International Medical Device Regulators Forum} (IMDRF). Entretanto, o gêmeo digital transita por múltiplas classes de risco simultaneamente: ele armazena dados de forma passiva (Classe I), processa imagens para auxílio diagnóstico (Classe II) e pode utilizar IA para sugerir opções invasivas, como posições críticas de implantes (Classe III ou IV). Esta fluidez taxonômica cria uma imensa insegurança jurídica. O desenvolvedor arca com custos milionários de submissão sem garantias sobre quais ensaios serão exigidos. Como aponta a extensa revisão sobre aspectos éticos e regulatórios \cite{springer2026realising}, o vácuo de diretrizes atua como o principal gargalo para investimentos da iniciativa privada.

\subsection{A Ausência de Padrões Mínimos de Calibração}

O mercado atual opera como um "velho oeste" digital. Cada fornecedor ou centro de pesquisa utiliza parâmetros idiossincráticos. Não existem normativas da ABNT ou de agências globais que delimitem a tolerância de erro geométrico máximo para a sobreposição entre IOS e CBCT (o \textit{registration error}). Não há padrões para a calibração de impressoras 3D utilizadas na etapa final, tampouco requisitos mandatórios de auditoria algorítmica para detectar vieses.

\destaque{A ausência de protocolos padronizados de validação significa que os resultados gerados por gêmeos digitais de diferentes plataformas são, na prática, cientificamente incomparáveis.}

Ainda que a recém-publicada norma brasileira ABNT NBR ISO/IEC 3173 \cite{abinc2026brasil} tenha instituído o vocabulário base, a tradução desses conceitos de alto nível para especificações operacionais aplicáveis à prática endodôntica, ortodôntica ou cirúrgica continua sendo uma lacuna abissal que aguarda a manifestação coordenada do Conselho Federal de Odontologia (CFO) e das sociedades científicas especializadas.
